Research on reproducing the canonical stylized facts of financial markets spans parametric return models, latent-state frameworks, and deep generative approaches.
This section reviews the key developments that motivate the present work.

\subsection{Stylized Facts and Parametric Models}

The Black-Scholes-Merton framework \citep{black1973pricing} assumes log-returns are i.i.d.\ Gaussian, but this conflicts with well-documented empirical regularities.
\citet{mandelbrot1963variation} demonstrated that speculative price changes exhibit heavy tails (leptokurtosis), \citet{fama1970efficient} confirmed negligible linear autocorrelation in raw returns, and \citet{schwert1989does} showed that the autocorrelation of absolute returns decays slowly---the signature of volatility clustering.
\citet{cont2001empirical} provides a systematic review confirming these regularities across asset classes and sampling frequencies.

GARCH models \citep{engle1982autoregressive, bollerslev1986generalized} capture conditional variance dynamics but require fat-tailed noise distributions to produce heavy tails and cannot represent discrete market regimes.
Stochastic volatility models \citep{heston1993closed} treat volatility as a latent process but require computationally intensive estimation.
Jump-diffusion models \citep{merton1976option, kou2002jump} add Poisson-driven jumps to the price process, producing heavy-tailed marginals, but lack a mechanism for volatility persistence: post-jump volatility returns immediately to its baseline level.

\subsection{Hidden Markov Models and Regime-Switching}

\citet{hamilton1989new} introduced the regime-switching HMM to economics, showing that U.S.\ GNP growth is well described by a two-state model.
HMMs have since become a standard tool for latent-state modeling of financial time series \citep{kim1999has, nguyen2018hidden, rossi2006volatility}.
\citet{ryden1998stylized} systematically evaluated HMMs against stylized facts: Gaussian-emission HMMs with few states reproduce heavy tails and the absence of return autocorrelation, but fail to capture volatility clustering because the ACF of absolute returns decays too quickly.
\citet{bulla2006stylized} showed that hidden semi-Markov models (HSMMs) improve volatility-clustering reproduction by allowing flexible sojourn-time distributions.

The discrete HMM framework of \citet{alswaidan2026hybrid} addressed the volatility-clustering limitation by adding a Poisson jump-duration mechanism that forced the model to dwell in tail states for empirically realistic durations.
That framework used Laplace quantile binning and direct frequency counting, bypassing the Baum-Welch algorithm entirely.
The present work replaces the discretization step with continuous Gaussian emissions learned via Baum-Welch, preserving the jump mechanism while eliminating quantization error, and extends the framework to volatility indices (VIX) whose statistical properties---positive skewness, mean reversion, extreme kurtosis---differ markedly from equity returns.

\subsection{The Baum-Welch Algorithm}

The Baum-Welch algorithm \citep{baum1970maximization} is the standard EM procedure for HMM parameter estimation.
It alternates between an E-step (computing posterior state probabilities via the forward-backward algorithm) and an M-step (re-estimating transition probabilities, emission parameters, and the initial state distribution).
\citet{bilmes1998gentle} provides a tutorial treatment.
A well-known limitation is sensitivity to initialization: poor starting values can lead to degenerate local optima where all states collapse to the global mean \citep{rabiner1986introduction}.
We address this with a quantile-based initialization strategy that seeds each state's emission parameters from a distinct segment of the sorted observation sequence.

\subsection{Deep Generative Models}

Deep generative models have attracted interest for financial time series generation.
\citet{takahashi2019modeling} found that GAN architectures produced sequences that visually resembled financial data but failed to reproduce ACF structure.
\citet{kwon2024can} confirmed that volatility clustering remained elusive unless the architecture was specifically designed for it.
\citet{alswaidan2026hybrid} benchmarked a GRU-based generator against HMM variants and found that it closed 83\% of the ACF-MAE gap but suffered variance collapse (0.6\% KS pass rate).
The present study focuses on the parametric HMM approach, which offers interpretability and computational efficiency advantages over neural generators.
